\textbf{Phần 1. Trắc nghiệm lựa chọn}
\setcounterpageref{ex}{0}

\begin{ex} %[3H1V07]
Một sợi dây đồng dài $2\text{ m } 35\text{ cm}$. Một sợi dây nhôm dài $24\text{ dm}$. Hỏi sợi dây nào dài hơn và dài hơn bao nhiêu xăng-ti-mét?
\choice
{Sợi dây đồng dài hơn $5\text{ cm}$}
{Sợi dây nhôm dài hơn $15\text{ cm}$}
{\True Sợi dây nhôm dài hơn $5\text{ cm}$}
{Sợi dây đồng dài hơn $15\text{ cm}$}
\loigiai{
Đổi các số đo về cùng đơn vị xăng-ti-mét:\\
- Sợi dây đồng dài: $2\text{ m } 35\text{ cm} = 235\text{ cm}$\\
- Sợi dây nhôm dài: $24\text{ dm} = 240\text{ cm}$\\
Vì $240\text{ cm} > 235\text{ cm}$ nên sợi dây nhôm dài hơn và dài hơn số xăng-ti-mét là:
$$240 - 235 = 5\text{ (cm)}$$
\textbf{Đáp án của câu hỏi là C}
}
\end{ex}

\begin{ex} %[3H1V07]
Đường gấp khúc $ABCD$ gồm 3 đoạn thẳng: đoạn thẳng $AB$ dài $18\text{ cm}$, đoạn thẳng $BC$ dài hơn đoạn thẳng $AB$ là $6\text{ cm}$, đoạn thẳng $CD$ dài bằng một nửa đoạn thẳng $BC$. Độ dài đường gấp khúc $ABCD$ là
\choice
{\True $54\text{ cm}$}
{$48\text{ cm}$}
{$56\text{ cm}$}
{$50\text{ cm}$}
\loigiai{
Độ dài đoạn thẳng $BC$ là:
$$18 + 6 = 24\text{ (cm)}$$
Độ dài đoạn thẳng $CD$ là:
$$24 : 2 = 12\text{ (cm)}$$
Độ dài đường gấp khúc $ABCD$ là:
$$18 + 24 + 12 = 54\text{ (cm)}$$
\textbf{Đáp án của câu hỏi là A}
}
\end{ex}

\begin{ex} %[3H1V07]
Một hình tam giác có độ dài cạnh thứ nhất là $15\text{ cm}$, cạnh thứ hai dài $20\text{ cm}$. Biết chu vi của hình tam giác đó là $52\text{ cm}$. Độ dài cạnh thứ ba của hình tam giác là
\choice
{$15\text{ cm}$}
{$16\text{ cm}$}
{$18\text{ cm}$}
{\True $17\text{ cm}$}
\loigiai{
Tổng độ dài hai cạnh thứ nhất và thứ hai là:
$$15 + 20 = 35\text{ (cm)}$$
Độ dài cạnh thứ ba của hình tam giác là:
$$52 - 35 = 17\text{ (cm)}$$
\textbf{Đáp án của câu hỏi là D}
}
\end{ex}

\begin{ex} %[3H1V07]
Trên đĩa cân bên trái đặt một quả mít và một quả cân $1\text{ kg}$. Trên đĩa cân bên phải đặt hai quả cân loại $2\text{ kg}$ và một quả cân $1\text{ kg}$ thì cân thăng bằng. Quả mít đó nặng số ki-lô-gam là
\choice
{$3\text{ kg}$}
{\True $4\text{ kg}$}
{$5\text{ kg}$}
{$2\text{ kg}$}
\loigiai{
Tổng khối lượng các quả cân trên đĩa cân bên phải là:
$$2 + 2 + 1 = 5\text{ (kg)}$$
Vì cân ở vị trí thăng bằng nên quả mít nặng số ki-lô-gam là:
$$5 - 1 = 4\text{ (kg)}$$
\textbf{Đáp án của câu hỏi là B}
}
\end{ex}

\begin{ex} %[3H1V07]
Có hai thùng chứa nước mắm. Thùng thứ nhất chứa $48\text{ lít}$ nước mắm. Nếu chuyển $8\text{ lít}$ nước mắm từ thùng thứ nhất sang thùng thứ hai thì hai thùng có số lít nước mắm bằng nhau. Lúc đầu thùng thứ hai chứa số lít nước mắm là
\choice
{$36\text{ lít}$}
{$40\text{ lít}$}
{\True $32\text{ lít}$}
{$24\text{ lít}$}
\loigiai{
Sau khi chuyển $8\text{ lít}$ nước mắm, thùng thứ nhất còn lại số lít là:
$$48 - 8 = 40\text{ (lít)}$$
Lúc này thùng thứ hai cũng có $40\text{ lít}$ nước mắm.\\
Ban đầu thùng thứ hai chứa số lít nước mắm là:
$$40 - 8 = 32\text{ (lít)}$$
\textbf{Đáp án của câu hỏi là C}
}
\end{ex}

\begin{ex} %[3H1V07]
Một mảnh vườn hình tứ giác có độ dài các cạnh lần lượt là $12\text{ m}, 15\text{ m}, 18\text{ m}$ và $25\text{ m}$. Người ta làm hàng rào bao quanh mảnh vườn đó và để lại một cổng vào rộng $2\text{ m}$. Độ dài hàng rào bao quanh mảnh vườn là
\choice
{\True $68\text{ m}$}
{$70\text{ m}$}
{$72\text{ m}$}
{$66\text{ m}$}
\loigiai{
Chu vi mảnh vườn hình tứ giác là:
$$12 + 15 + 18 + 25 = 70\text{ (m)}$$
Độ dài hàng rào bao quanh mảnh vườn là:
$$70 - 2 = 68\text{ (m)}$$
\textbf{Đáp án của câu hỏi là A}
}
\end{ex}

\begin{ex} %[3H1V07]
Tính giá trị của biểu thức $1\text{ m } 8\text{ dm} - 45\text{ cm} + 2\text{ dm}$ khi đổi ra xăng-ti-mét:
\choice
{$145\text{ cm}$}
{$165\text{ cm}$}
{$150\text{ cm}$}
{\True $155\text{ cm}$}
\loigiai{
Đổi các số đo sang đơn vị xăng-ti-mét:\\
- $1\text{ m } 8\text{ dm} = 180\text{ cm}$\\
- $2\text{ dm} = 20\text{ cm}$\\
Thực hiện phép tính:\\
$$180\text{ cm} - 45\text{ cm} + 20\text{ cm} = 135\text{ cm} + 20\text{ cm} = 155\text{ cm}$$
\textbf{Đáp án của câu hỏi là D}
}
\end{ex}

\begin{ex} %[3H1V07]
Đoạn thẳng $AB$ dài $16\text{ cm}$. Điểm $M$ là trung điểm của đoạn thẳng $AB$. Điểm $N$ là trung điểm của đoạn thẳng $AM$. Độ dài đoạn thẳng $AN$ là
\choice
{$2\text{ cm}$}
{\True $4\text{ cm}$}
{$6\text{ cm}$}
{$8\text{ cm}$}
\loigiai{
Vì $M$ là trung điểm của đoạn thẳng $AB$ nên độ dài đoạn thẳng $AM$ là:
$$16 : 2 = 8\text{ (cm)}$$
Vì $N$ là trung điểm của đoạn thẳng $AM$ nên độ dài đoạn thẳng $AN$ là:
$$8 : 2 = 4\text{ (cm)}$$
\textbf{Đáp án của câu hỏi là B}
}
\end{ex}

\begin{ex} %[3H1V07]
Một bao gạo cân nặng $50\text{ kg}$. Người ta lấy ra một nửa số gạo trong bao rồi đổ thêm vào bao $5\text{ kg}$ gạo nếp. Lúc này trong bao có tất cả số ki-lô-gam gạo là
\choice
{$25\text{ kg}$}
{$35\text{ kg}$}
{\True $30\text{ kg}$}
{$28\text{ kg}$}
\loigiai{
Một nửa số gạo trong bao là:
$$50 : 2 = 25\text{ (kg)}$$
Lúc này trong bao có tất cả số ki-lô-gam gạo là:
$$25 + 5 = 30\text{ (kg)}$$
\textbf{Đáp án của câu hỏi là C}
}
\end{ex}

\begin{ex} %[3H1V07]
Một bình chứa $1\text{ lít}$ nước. Người ta rót nước từ bình vào đầy $3$ chiếc cốc loại $200\text{ ml}$ và $1$ chiếc cốc loại $250\text{ ml}$. Lượng nước còn lại trong bình là
\choice
{\True $150\text{ ml}$}
{$200\text{ ml}$}
{$100\text{ ml}$}
{$250\text{ ml}$}
\loigiai{
Đổi: $1\text{ lít} = 1000\text{ ml}$.\\
Lượng nước đã rót vào các cốc là:
$$3 \times 200 + 250 = 600 + 250 = 850\text{ (ml)}$$
Lượng nước còn lại trong bình là:
$$1000 - 850 = 150\text{ (ml)}$$
\textbf{Đáp án của câu hỏi là A}
}
\end{ex}

\begin{ex} %[3H1V07]
Một hình tam giác có độ dài ba cạnh bằng nhau, biết chu vi của hình tam giác đó bằng $24\text{ cm}$. Độ dài mỗi cạnh của hình tam giác đó là
\choice
{$6\text{ cm}$}
{$7\text{ cm}$}
{$9\text{ cm}$}
{\True $8\text{ cm}$}
\loigiai{
Vì hình tam giác có ba cạnh bằng nhau nên độ dài mỗi cạnh là:
$$24 : 3 = 8\text{ (cm)}$$
\textbf{Đáp án của câu hỏi là D}
}
\end{ex}

\begin{ex} %[3H1V07]
Con gà cân nặng $2\text{ kg}$, con vịt nặng hơn con gà $1\text{ kg}$. Con ngỗng cân nặng gấp đôi con vịt. Hỏi con ngỗng nặng hơn con gà bao nhiêu ki-lô-gam?
\choice
{$3\text{ kg}$}
{\True $4\text{ kg}$}
{$5\text{ kg}$}
{$2\text{ kg}$}
\loigiai{
Cân nặng của con vịt là:
$$2 + 1 = 3\text{ (kg)}$$
Cân nặng của con ngỗng là:
$$3 \times 2 = 6\text{ (kg)}$$
Con ngỗng nặng hơn con gà số ki-lô-gam là:
$$6 - 2 = 4\text{ (kg)}$$
\textbf{Đáp án của câu hỏi là B}
}
\end{ex}

\textbf{Phần 2. Trắc nghiệm đúng sai}
\setcounterpageref{ex}{0}

\begin{ex} %[3H1V07]
Xét tính đúng sai của các nhận định sau về đơn vị đo và hình học:
\choiceTF
{\True Mối quan hệ giữa các đơn vị đo độ dài: $1\text{ m} = 10\text{ dm} = 100\text{ cm}$.}
{$1\text{ kg}$ bông nhẹ hơn $1\text{ kg}$ sắt.}
{\True Nửa lít nước tương ứng với $500\text{ ml}$ nước.}
{\True Chu vi hình tam giác có ba cạnh dài $5\text{ cm}, 6\text{ cm}, 7\text{ cm}$ bằng chu vi hình chữ nhật có chiều dài $5\text{ cm}$ và chiều rộng $4\text{ cm}$.}
\loigiai{
\begin{itemize}
    \item Ý a) Đúng: Theo bảng đơn vị đo độ dài.
    \item Ý b) Sai: Cả hai đều nặng 1 kg nên có khối lượng bằng nhau.
    \item Ý c) Đúng: $1\text{ lít} = 1000\text{ ml}$, nửa lít là $1000 : 2 = 500\text{ ml}$.
    \item Ý d) Đúng: Chu vi tam giác: $5 + 6 + 7 = 18\text{ cm}$. Chu vi hình chữ nhật: $(5 + 4) \times 2 = 18\text{ cm}$. Hai chu vi bằng nhau.
\end{itemize}
\begin{itemize}
    \item a) \textbf{Đúng}.
    \item b) \textbf{Sai}.
    \item c) \textbf{Đúng}.
    \item d) \textbf{Đúng}.
\end{itemize}
}
\end{ex}

\begin{ex} %[3H1V07]
Cho đoạn thẳng $AB = 24\text{ cm}$. Lấy điểm $M$ là trung điểm của đoạn thẳng $AB$, sau đó lấy điểm $I$ là trung điểm của đoạn thẳng $AM$:
\choiceTF
{\True Độ dài đoạn thẳng $AM$ bằng $12\text{ cm}$.}
{\True Điểm $M$ nằm giữa hai điểm $A$ và $B$.}
{Độ dài đoạn thẳng $AI$ bằng $8\text{ cm}$.}
{\True Độ dài đoạn thẳng $IB$ bằng $18\text{ cm}$.}
\loigiai{
\begin{itemize}
    \item Ý a) Đúng: $M$ là trung điểm $AB$ nên $AM = 24 : 2 = 12\text{ (cm)}$.
    \item Ý b) Đúng: Trung điểm luôn nằm giữa hai đầu mút của đoạn thẳng.
    \item Ý c) Sai: $I$ là trung điểm $AM$ nên $AI = 12 : 2 = 6\text{ (cm)}$.
    \item Ý d) Đúng: Độ dài đoạn thẳng $IB = IM + MB = 6 + 12 = 18\text{ (cm)}$.
\end{itemize}
\begin{itemize}
    \item a) \textbf{Đúng}.
    \item b) \textbf{Đúng}.
    \item c) \textbf{Sai}.
    \item d) \textbf{Đúng}.
\end{itemize}
}
\end{ex}

\textbf{Phần 3. Trả lời ngắn}
\setcounterpageref{ex}{0}

\begin{ex} %[3H1V07]
Tính chu vi của một hình tam giác có độ dài các cạnh lần lượt là $14\text{ cm}, 18\text{ cm}$ và $23\text{ cm}$.
\shortans{55}
\loigiai{
\textbf{Đáp số:} $55$\\
Chu vi hình tam giác là:
$$14 + 18 + 23 = 55\text{ (cm)}$$
}
\end{ex}

\begin{ex} %[3H1V07]
Một sợi dây thép dài $3\text{ m}$ được uốn thành một hình tam giác có ba cạnh bằng nhau. Hỏi mỗi cạnh của hình tam giác đó dài bao nhiêu đề-xi-mét?
\shortans{10}
\loigiai{
\textbf{Đáp số:} $10$\\
Đổi: $3\text{ m} = 30\text{ dm}$.\\
Mỗi cạnh của hình tam giác dài số đề-xi-mét là:
$$30 : 3 = 10\text{ (dm)}$$
}
\end{ex}

\begin{ex} %[3H1V07]
Một con ngỗng cân nặng $6\text{ kg}$, một con gà cân nặng $2\text{ kg}$. Hỏi con ngỗng nặng gấp mấy lần con gà?
\shortans{3}
\loigiai{
\textbf{Đáp số:} $3$\\
Con ngỗng nặng gấp con gà số lần là:
$$6 : 2 = 3\text{ (lần)}$$
}
\end{ex}

\begin{ex} %[3H1V07]
Thùng thứ nhất chứa $35\text{ lít}$ dầu, thùng thứ hai chứa số lít dầu gấp $2$ lần thùng thứ nhất. Hỏi cả hai thùng chứa tất cả bao nhiêu lít dầu?
\shortans{105}
\loigiai{
\textbf{Đáp số:} $105$\\
Thùng thứ hai chứa số lít dầu là:
$$35 \times 2 = 70\text{ (lít)}$$
Cả hai thùng chứa tất cả số lít dầu là:
$$35 + 70 = 105\text{ (lít)}$$
}
\end{ex}

\begin{ex} %[3H1V07]
Đoạn thẳng $PQ$ dài $32\text{ cm}$. Điểm $K$ là trung điểm của đoạn thẳng $PQ$. Tính độ dài đoạn thẳng $PK$.
\shortans{16}
\loigiai{
\textbf{Đáp số:} $16$\\
Độ dài đoạn thẳng $PK$ là:
$$32 : 2 = 16\text{ (cm)}$$
}
\end{ex}

\begin{ex} %[3H1V07]
Người ta cắm $5$ chiếc cọc thẳng hàng cách đều nhau, khoảng cách giữa hai chiếc cọc liền kề là $3\text{ m}$. Hỏi khoảng cách từ chiếc cọc đầu tiên đến chiếc cọc cuối cùng là bao nhiêu mét?
\shortans{12}
\loigiai{
\textbf{Đáp số:} $12$\\
5 chiếc cọc cắm thẳng hàng chia đoạn đường thành: $5 - 1 = 4$ (khoảng cách).\\
Khoảng cách từ chiếc cọc đầu tiên đến chiếc cọc cuối cùng là:
$$4 \times 3 = 12\text{ (m)}$$
}
\end{ex}
